% GENERATED FILE. Edit canonical lessons, then run npm run generate:books.
% canonical-source-hash: fnv1a64:a75f923a5535efdd
% canonical-lessons: GE-C469-enttaeuschen, GE-C469-enttaeuschung, GE-C469-erleichterung, GE-C469-aufatmen, GE-C469-seufzen

\chapter{Enttäuscht und erleichtert}
\label{ch:ge-enttauscht-und-erleichtert}
\hlchaptermodality{469}

\begin{chapteropening}
\textbf{By the end of this chapter:} \emph{I can talk about disappointment and relief.}

Complete the last lesson of chapter 469: I can talk about disappointment and relief.
\end{chapteropening}

\section[enttäuschen]{enttäuschen — to disappoint}

\label{lesson:GE-C469-enttaeuschen}



\begin{quote}
Before the new one: say the German for to get worked up, then the German for to rage.
\end{quote}

\subsection*{You'll want to know: enttäuschen}
\textbf{enttäuschen} — \textquotedblleft{}to disappoint\textquotedblright{}.

\begin{quote}
Der Film hat mich enttäuscht.
\end{quote}

\begin{quote}
\emph{The film disappointed me.}
\end{quote}

\emph{Enttäuscht sein von} — to be disappointed by or in. Inseparable: \emph{hat enttäuscht}.

\begin{cousinweb}
\emph{enttäuschen} is built on \emph{ent-} \textquotedblleft{}away\textquotedblright{} and \emph{täuschen} \textquotedblleft{}to deceive\textquotedblright{}, from Middle High German \emph{tiuschen} (EtymDB); it is usually explained as first meaning \textquotedblleft{}to free someone from a delusion\textquotedblright{}.
\end{cousinweb}

\begin{grammarlens}[title={What you've built}]
One more word for telling, explaining and describing at B1.
\end{grammarlens}

\subsection*{Your turn}
\begin{itemize}
\raggedright
  \item \emph{Say it:} \emph{enttäuschen}
  \item \emph{Say it:} \emph{enttäuschen}, once more
  \item \emph{Say it:} a sentence of your own with \emph{enttäuschen}
  \item \emph{Recall:} say the German for to get worked up, then the German for to rage, then say \emph{enttäuschen} again
\end{itemize}

\begin{tcolorbox}[breakable,colback=teal!4,colframe=teal!35!black,title={Before you move on}]
What does enttäuschen mean? (\textquotedblleft{}To disappoint\textquotedblright{}.) Say it once more.
\end{tcolorbox}
\section[die Enttäuschung]{die Enttäuschung — disappointment}

\label{lesson:GE-C469-enttaeuschung}



\begin{quote}
Before the new one: say the German for to rage, then the German for to disappoint.
\end{quote}

\subsection*{You'll want to know: die Enttäuschung}
\textbf{die Enttäuschung} — \textquotedblleft{}disappointment\textquotedblright{}.

\begin{quote}
Die Niederlage war eine große Enttäuschung für die Fans.
\end{quote}

\begin{quote}
\emph{The defeat was a big disappointment for the fans.}
\end{quote}

Plural: \emph{die Enttäuschungen}.

\begin{cousinweb}
\emph{Enttäuschung} is built on \emph{enttäuschen} with \emph{-ung}; the verb joins \emph{ent-} \textquotedblleft{}away\textquotedblright{} to \emph{täuschen} \textquotedblleft{}to deceive\textquotedblright{}, from Middle High German \emph{tiuschen} (EtymDB).
\end{cousinweb}

\begin{grammarlens}[title={What you've built}]
One more word for telling, explaining and describing at B1.
\end{grammarlens}

\subsection*{Your turn}
\begin{itemize}
\raggedright
  \item \emph{Say it:} \emph{die Enttäuschung}
  \item \emph{Say it:} \emph{die Enttäuschung}, once more
  \item \emph{Say it:} a sentence of your own with \emph{die Enttäuschung}
  \item \emph{Recall:} say the German for to rage, then the German for to disappoint, then say \emph{die Enttäuschung} again
\end{itemize}

\begin{tcolorbox}[breakable,colback=teal!4,colframe=teal!35!black,title={Before you move on}]
What does die Enttäuschung mean? (\textquotedblleft{}Disappointment\textquotedblright{}.) Say it once more.
\end{tcolorbox}
\section[die Erleichterung]{die Erleichterung — relief}

\label{lesson:GE-C469-erleichterung}



\begin{quote}
Before the new one: say the German for to disappoint, then the German for disappointment.
\end{quote}

\subsection*{You'll want to know: die Erleichterung}
\textbf{die Erleichterung} — \textquotedblleft{}relief\textquotedblright{}.

\begin{quote}
Zu unserer Erleichterung war niemand verletzt.
\end{quote}

\begin{quote}
\emph{To our relief, nobody was hurt.}
\end{quote}

No plural in this sense. \emph{Erleichtert} — relieved.

\begin{cousinweb}
\emph{Erleichterung} is built on \emph{erleichtern} \textquotedblleft{}to make lighter\textquotedblright{}, from \emph{leicht} \textquotedblleft{}light\textquotedblright{}, which is inherited from Old High German \emph{līht}, from Proto-Germanic \emph{linhtaz} (EtymDB).
\end{cousinweb}

\begin{grammarlens}[title={What you've built}]
One more word for telling, explaining and describing at B1.
\end{grammarlens}

\subsection*{Your turn}
\begin{itemize}
\raggedright
  \item \emph{Say it:} \emph{die Erleichterung}
  \item \emph{Say it:} \emph{die Erleichterung}, once more
  \item \emph{Say it:} a sentence of your own with \emph{die Erleichterung}
  \item \emph{Recall:} say the German for to disappoint, then the German for disappointment, then say \emph{die Erleichterung} again
\end{itemize}

\begin{tcolorbox}[breakable,colback=teal!4,colframe=teal!35!black,title={Before you move on}]
What does die Erleichterung mean? (\textquotedblleft{}Relief\textquotedblright{}.) Say it once more.
\end{tcolorbox}
\section[aufatmen]{aufatmen — to breathe a sigh of relief}

\label{lesson:GE-C469-aufatmen}



\begin{quote}
Before the new one: say the German for disappointment, then the German for relief.
\end{quote}

\subsection*{You'll want to know: aufatmen}
\textbf{aufatmen} — \textquotedblleft{}to breathe a sigh of relief\textquotedblright{}.

\begin{quote}
Nach der Prüfung atmeten alle auf.
\end{quote}

\begin{quote}
\emph{After the exam, everyone breathed a sigh of relief.}
\end{quote}

Separable: \emph{ich atme auf}, \emph{hat aufgeatmet}.

\begin{cousinweb}
\emph{aufatmen} is built on \emph{auf} \textquotedblleft{}up\textquotedblright{} and \emph{atmen} \textquotedblleft{}to breathe\textquotedblright{}, inherited from Old High German \emph{ātamōn}, from Proto-Germanic \emph{ēþmaz} \textquotedblleft{}breath\textquotedblright{} (EtymDB).
\end{cousinweb}

\begin{grammarlens}[title={What you've built}]
One more word for telling, explaining and describing at B1.
\end{grammarlens}

\subsection*{Your turn}
\begin{itemize}
\raggedright
  \item \emph{Say it:} \emph{aufatmen}
  \item \emph{Say it:} \emph{aufatmen}, once more
  \item \emph{Say it:} a sentence of your own with \emph{aufatmen}
  \item \emph{Recall:} say the German for disappointment, then the German for relief, then say \emph{aufatmen} again
\end{itemize}

\begin{tcolorbox}[breakable,colback=teal!4,colframe=teal!35!black,title={Before you move on}]
What does aufatmen mean? (\textquotedblleft{}To breathe a sigh of relief\textquotedblright{}.) Say it once more.
\end{tcolorbox}
\section[seufzen]{seufzen — to sigh}

\label{lesson:GE-C469-seufzen}



\begin{quote}
Before the new one: say the German for relief, then the German for to breathe a sigh of relief.
\end{quote}

\subsection*{You'll want to know: seufzen}
\textbf{seufzen} — \textquotedblleft{}to sigh\textquotedblright{}.

\begin{quote}
Sie sah auf die Uhr und seufzte.
\end{quote}

\begin{quote}
\emph{She looked at the clock and sighed.}
\end{quote}

Regular: \emph{hat geseufzt}. The noun: \emph{der Seufzer} — a sigh.

\begin{cousinweb}
\emph{seufzen} comes from Middle High German \emph{siufzen} (EtymDB); its older history is not settled.
\end{cousinweb}

\begin{grammarlens}[title={What you've built}]
One more word for telling, explaining and describing at B1.
\end{grammarlens}

\subsection*{Your turn}
\begin{itemize}
\raggedright
  \item \emph{Say it:} \emph{seufzen}
  \item \emph{Say it:} \emph{seufzen}, once more
  \item \emph{Say it:} a sentence of your own with \emph{seufzen}
  \item \emph{Recall:} say the German for relief, then the German for to breathe a sigh of relief, then say \emph{seufzen} again
\end{itemize}

\begin{tcolorbox}[breakable,colback=teal!4,colframe=teal!35!black,title={Before you move on}]
What does seufzen mean? (\textquotedblleft{}To sigh\textquotedblright{}.) Say it once more.
\end{tcolorbox}
